\documentclass{article}
\usepackage[utf8]{inputenc}
\usepackage[T1]{fontenc}
\usepackage{listings}
\usepackage{graphicx}
\usepackage{color}
\usepackage[top=15mm, bottom=6mm, left=13mm, right=13mm]{geometry}
\pagenumbering{gobble}
\usepackage{textgreek}

% Packages for headers and footers
\usepackage{fancyhdr}
\pagestyle{fancy}

% Define your custom R style
\definecolor{codegreen}{rgb}{0,0.6,0}
\definecolor{codegray}{rgb}{0.5,0.5,0.5}
\definecolor{codepurple}{rgb}{0.58,0,0.82}
\definecolor{backcolour}{rgb}{0.95,0.95,0.92}

\lstdefinestyle{mystyle}{
    backgroundcolor=\color{backcolour},   
    commentstyle=\color{codegreen},
    keywordstyle=\color{magenta},
    numberstyle=\tiny\color{codegray},
    stringstyle=\color{codepurple},
    basicstyle=\ttfamily\scriptsize,  % Change this line
    breakatwhitespace=false,         
    breaklines=true,                 
    keepspaces=true,                 
    numbers=left,                    
    numbersep=5pt,                  
    showspaces=false,                
    showstringspaces=false,
    showtabs=false,                  
    tabsize=2,
    literate=%
        {á}{{\'a}}1
        {é}{{\'e}}1
        {í}{{\'i}}1
        {ó}{{\'o}}1
        {ú}{{\'u}}1
        {Á}{{\'A}}1
        {É}{{\'E}}1
        {Í}{{\'I}}1
        {Ó}{{\'O}}1
        {Ú}{{\'U}}1
        {à}{{\`a}}1
        {è}{{\`e}}1
        {ì}{{\`i}}1
        {ò}{{\`o}}1
        {ù}{{\`u}}1
        {À}{{\`A}}1
        {È}{{\`E}}1
        {Ì}{{\`I}}1
        {Ò}{{\`O}}1
        {Ù}{{\`U}}1
        {ã}{{\~a}}1
        {õ}{{\~o}}1
        {Ã}{{\~A}}1
        {â}{{\~A}}1
        {Õ}{{\~O}}1
        {ç}{{\c c}}1
        {Ç}{{\c C}}1
        {α}{{\textalpha}}1
        {β}{{\textbeta}}1
        {μ}{{\textmu}}1
        {σ}{{\textsigma}}1
        {γ}{{\textgamma}}1
}

\lstset{style=mystyle}

% Define your headers
\rhead{Miguel Freitas (102756)}
\lhead{Afonso Ribeiro (102763)} % left side empty
\chead{Luís Correia (103528)} % center empty



\begin{document}
\centering
\noindent{\Large{\textbf{Projeto Computacional de Probabilidade e 
Estatística – Exercício 9}}}
\vspace{3mm}

\begin{lstlisting}[language=R, xleftmargin=4em, xrightmargin=4em]
set.seed(1593)
k <- 1000
p <- 0.5
γ <- 0.91
α <- 1 - γ
n_values <- c(30, 50, 100, 200, 300, 500, 1000)
m <- length(n_values)

medias_diferencas <- numeric(length(n_values))

for (i in 1:m) {
  diferencas <- numeric(k)
  n <- n_values[i]
  
  for (j in 1:k) {
    X <- rbinom(n, size = 1, prob = p)
    
    # Método 1
    media <- mean(X)
    z <- qnorm((1 + γ)/2, mean = 0, sd = 1)
    
      ## Resolução da equação de segundo grau:
      a <- (1 + z^2 / n)
      b <- -2 * media - z^2 / n
      c <- media^2
      solucao1 <- (-b - sqrt(b^2 - 4*a*c)) / (2*a)
      solucao2 <- (-b + sqrt(b^2 - 4*a*c)) / (2*a)
      comprimento1 <- solucao2 - solucao1
    
    # Método 2
    b_ <- qnorm(1 - α/2, mean = 0, sd = 1)
    limite_inferior <- media - b_ * sqrt( (media * (1 - media)) / n_values[i])
    limite_superior <- media + b_ * sqrt( (media * (1 - media)) / n_values[i])
    comprimento2 <- limite_superior - limite_inferior
    
    # Diferença entre os comprimentos
    diferencas[j] <- comprimento2 - comprimento1
  }
  
  medias_diferencas[i] <- mean(diferencas)
}

# Criar o gráfico
dev.new()
par(mar = c(5, 5, 4, 2))
plot(n_values, medias_diferencas, col = "blue", pch = 19, type = "b",
     xlab = "Tamanho da amostra",
     ylab = "Diferença entre os comprimentos dos intervalos de confiança",
     cex.lab = 1.4, cex.main = 1.5, cex.axis = 1.3, cex = 1.5,
     main = paste("Diferenças médias entre os comprimentos dos intervalos de confiança\n",
                  "construídos pelo Método 2 e pelo Método 1, em função do tamanho da amostra"))
\end{lstlisting}

\vspace{0.1cm}

\begin{figure}[ht]
\begin{minipage}{0.63\textwidth}
\centering
\includegraphics[width=\linewidth, height=0.7\textheight, keepaspectratio]{ex9.jpeg}
\label{fig:mylabel}
\end{minipage}
\begin{minipage}{0.3\textwidth}
Para encontrar os limites do intervalo de confiança do parâmetro p, deve-se isolar o valor de p e encontrar as raízes do polinómio de 2.º grau, tal como é explicado no método 1. O método 2 baseia-se no facto de que $\bar{X}$ é um estimador consistente de p e, como tal, substitui-se p no denominador de Z por $\bar{X}$, o que facilita os cálculos.
Tendo isto em conta, é natural que, quanto maior o tamanho da amostra, melhor é a estimação de p por $\bar{X}$, ou seja, menor será o erro da utilização do estimador de p ao invés de p, no método 2. Isto é corroborado pelo gráfico obtido, que mostra que, com o aumento do tamanho da amostra, a diferença entre os comprimentos dos intervalos de confiança construídos pelo método 2 e pelo método 1 diminui, ou seja, vão-se tornando cada vez mais equiparáveis.
\end{minipage}
\end{figure}

\end{document}
